%---------------------------------------------------------------
%  SECTION IV — RESULTS AND EVALUATION
%---------------------------------------------------------------
\section{Results and Evaluation}
\label{sec:results}

\subsection{Handshake Validation}

The Q-TLS handshake was tested in both localhost and cross-device
configurations. Figures~\ref{fig:client_out} and~\ref{fig:server_out}
show the CLI output confirming successful certificate verification,
Kyber512 key encapsulation, Dilithium2 transcript signing, and
HKDF-SHA256 session key derivation.

\begin{figure}[htbp]
  \centering
  \includegraphics[width=\columnwidth]{client}
  \caption{Client-side Q-TLS handshake output (all four phases).}
  \label{fig:client_out}
\end{figure}

\begin{figure}[htbp]
  \centering
  \includegraphics[width=\columnwidth]{server}
  \caption{Server-side Q-TLS handshake output and session key establishment.}
  \label{fig:server_out}
\end{figure}

\subsection{Cryptographic Parameters}

Table~\ref{tab:crypto} summarises the algorithm parameters.
All key sizes are as specified in FIPS~203/204 \cite{nist2024pqc}.

\begin{table}[htbp]
  \centering
  \caption{Cryptographic Parameters in Q-SFTP}
  \label{tab:crypto}
  \begin{tabular}{lcccc}
    \toprule
    \textbf{Algorithm} & \textbf{Level} & \textbf{PK} & \textbf{SK}
                       & \textbf{CT / Sig} \\
    \midrule
    Kyber512     & NIST L1        & 800 B   & 1,632 B & 768 B \\
    Dilithium2   & NIST L2        & 1,312 B & 2,560 B & 2,420 B \\
    AES-256-GCM  & 128-bit (PQ)   & ---     & 32 B    & +28 B$^{\dag}$ \\
    BLAKE2b-256  & 128-bit (PQ)   & ---     & ---     & 32 B \\
    \bottomrule
    \multicolumn{5}{l}{$^{\dag}$12 B nonce + 16 B authentication tag per message.}
  \end{tabular}
\end{table}

\subsection{Performance Benchmarks}

Table~\ref{tab:dilithium} shows Dilithium performance benchmarks
and Table~\ref{tab:kyber} shows Kyber benchmarks, both measured on an
Intel Core i7-9750H CPU averaged over 1000 iterations.

\begin{table}[htbp]
  \centering
  \caption{Dilithium Digital Signature Benchmarks}
  \label{tab:dilithium}
  \begin{tabular}{lccc}
    \toprule
    \textbf{Operation} & \textbf{Dil.2} & \textbf{Dil.3}
                       & \textbf{Dil.5} \\
    \midrule
    KeyGen()  & 6 ms  & 9 ms  & 15 ms \\
    Sign()    & 27 ms & 46 ms & 58 ms \\
    Verify()  & 7 ms  & 11 ms & 18 ms \\
    \bottomrule
  \end{tabular}
\end{table}

\begin{table}[htbp]
  \centering
  \caption{Kyber KEM Benchmarks (time / throughput)}
  \label{tab:kyber}
  \begin{tabular}{lcccccc}
    \toprule
    \textbf{Params} & \textbf{KG} & \textbf{KG/s}
                    & \textbf{Enc} & \textbf{Enc/s}
                    & \textbf{Dec} & \textbf{Dec/s} \\
    \midrule
    Kyber512  & 2.0 ms & 494 & 2.8 ms & 353 & 4.1 ms & 243 \\
    Kyber768  & 3.4 ms & 294 & 4.4 ms & 228 & 6.1 ms & 165 \\
    Kyber1024 & 5.1 ms & 197 & 6.2 ms & 161 & 8.3 ms & 121 \\
    \bottomrule
  \end{tabular}
\end{table}

Q-SFTP uses Kyber512, completing a full handshake (two Dilithium
operations + one Kyber encaps/decaps) in under 80~ms on the same
hardware, which is acceptable for a per-session operation.

\subsection{Chunked Transfer Overhead}

For a file of size $F$ with 256~KB chunks:
$N = \lceil F / 262\,144 \rceil$ chunks are transmitted.
Each chunk incurs 28 bytes of AES-256-GCM overhead. For a 100~MB
file ($N = 400$ chunks), the total overhead is $400 \times 28
= 11{,}200$ bytes, or approximately $\mathbf{0.01\%}$.

\subsection{Automated Security Testing}

\begin{figure}[htbp]
  \centering
  \includegraphics[width=\columnwidth]{pytest_handshake}
  \caption{Pytest test suite results for the Q-TLS handshake.}
  \label{fig:pytest_hs}
\end{figure}

\begin{figure}[htbp]
  \centering
  \includegraphics[width=\columnwidth]{pytest_ca}
  \caption{Pytest test suite results for the CA module.}
  \label{fig:pytest_ca}
\end{figure}

As shown in Figs.~\ref{fig:pytest_hs} and \ref{fig:pytest_ca},
100\% of test cases pass across certificate issuance, mutual
authentication, key derivation, hash verification, and replay
detection. The security properties verified are:

\begin{itemize}
  \item \textbf{Quantum resistance}: IND-CCA2 under Kyber512 and
        EUF-CMA under Dilithium2 (MLWE/MSIS hardness assumptions).
  \item \textbf{MITM prevention}: Bilateral Dilithium2 transcript
        signatures prevent impersonation.
  \item \textbf{Replay prevention}: Random nonces and 30-second
        timestamp window.
  \item \textbf{Data integrity}: AES-256-GCM tags (in-transit) and
        BLAKE2b-256 hashes (at-rest).
  \item \textbf{Privacy}: Metadata stripped before storage; IPs
        anonymised via rotating SHA-256 salts.
\end{itemize}
